\section{The Unum Necessarium}

\emph{``Papers, comrade!''}


Many traditionalists are going to argue that initiation must transpire through the conduit or under the aegis of an ``Order''. That this is a desirable state of affairs is (of course) obvious. When Voldemort's power is exalted, one looks around for an ``Order of the Phoenix'' to resist him. This situation, otherwise normative, is obviated during the ``end of days'', since the very definition of the Kali Yuga is that of an era devoid of such influence, an epoch where the ``powers of evil are exalted'' during the time of the mystery of iniquity.

No one is arguing that (perhaps) someone may stumble upon a genuine initiation from a master, or that they are wrong to take advantage of such. However, the vast majority of those who choose to become part of the remnant are going to have little to no access to such aid. ``The harvest is great, but the workers are few''. Nevertheless, to stand in such a day requires that one ``ride the tiger'', which in this case, means \emph{self-initiation}.

That which would otherwise be forbidden, and in an earlier time, a sign of great foolishness, becomes \emph{the one thing necessary}. Just as many Christians are going to misunderstand this path as self-deification, so many pagans (it would seem) are going to take umbrage at this path as self-teaching or self-glorification. All this proves is that both parties (which hate the other) actually think pretty much alike. They have a deep, dogmatic disdain for anyone who is unwilling to sit around \& wait for non-existent superiors to take action to save the sinking Ark. If they could, they would anathematize those who are willing to do the wrong thing (if need be) rather than risk doing nothing.

Time is not a luxury which any human possesses at any period of time. For us, far less so. The ``Norman'' class of those who are willing to fundamentally allow business as usual to continue are not going to like it when the Hooded Man\footnote{\url{https://www.youtube.com/watch?v=OnkWvh1pvU8}} returns to Sherwood Forest. Left to their tender mercies, the path of self-initiation is going to be either endlessly ``tabled'', or forever ``disqualified'' due to technicalities.

The alchemist will not be dismayed at this; in fact, he will take fuel and flame from the ``shock'' of how erstwhile brethren are actually ``enemies'' and ``foes'' in this path. It is all part of this path, this strange way of the Self. ``Strike, Arjuna\footnote{\url{https://en.wikipedia.org/wiki/Arjuna}},'' bade the god at the end of time, ``and fear not.''

For the blow is against, and for, the Self.

Many of the problems with traditionalism in the 20th century seem to have arisen because various schools were ``lukewarm'' -- that is, they were unable to transcend themselves enough to overcome various difficulties along the path. This manifested in in-fighting \& schisms
\& fruitless disagreements.


There are exceptions to this. Peter Deunov\footnote{\url{https://en.wikipedia.org/wiki/Peter\_Deunov}} established contact with Krishmanurti\footnote{\url{https://gornahoor.net/?p=2457}} during the time of the dissolution of the Order of the Eastern Star. Here, again, we see how the Slavic soul has been able to evade so many of the problems which overtook the Western soul; Theosophy and Anthroposophy in the West became hopelessly muddled and mired over personalities and issues, so much so that not even Steiner's genius and influence could entirely evade the trend.

Men like Deunov, Gurdjieff, Bondarev, and many others (Tomasso Palamidessi and Franz Bardon were not part of the Anglo-sphere) were able to a great extent to steer clear of the Theosophy fiasco. Why is this?

Even Evola evaded the dangers of compromise and dilution, but arguably failed to rise to his own occasion, and we have seen how Guenon thought that conversion to Islam\footnote{\url{http://theforbiddenreligion.com/recursos/the-forbidden-religion.pdf}} was the answer to the West's bankruptcy:

\begin{quotex}
It took Rene Guenon, for example, years of study and meditation to make the decision to abandon christianity and join masonry and martinism, only to give it all up at a later date and convert to islam. He thought he had made great leaps with these changes, but the only thing he did was to go round in circles inside his labyrinth. His search was in vain. And if Guenon, well-learned in these subjects, experienced such confusion, one can imagine the mistakes the average man will make.
\end{quotex}


This may not be a totally accurate assessment, but it is critical to build upon the past, while avoiding certain turns or twists or dead ends. To do this, we must acquit ourselves as some of the better among them have -- like Arjuna, who could only see the eye of the bird that he was asked to hit with the bow, they were focused on the one thing necessary.

It is this which makes Gornahoor's task difficult. It would be easy to hit the bird, and possibly not too difficult to hit the eye. But to merely see the eye, and to exclude all else -- that is not possible for any but those who are chosen. Hence, the many obstacles and difficulties -- it would be easy, for instance, to not be constrained by the patrimony of the Middle Ages. Men like Deunov and Krishmanurti demonstrate what is possible (beyond men like Steiner) when there is a desire to go farther.

Romanity, whether that of Rome itself, or the Third Rome, is an idea that has helped clarify and distill much esoteric practice. The Slavic soul has not lost, nor will it lose, the remembrance of ``the glory that was Greece, the grandeur that was Rome\footnote{\url{https://gornahoor.net/?p=223}}``.

Romanity is like the guardian of the lower threshhold\footnote{\url{https://en.wikipedia.org/wiki/Guardian\_of\_the\_Threshold}}, the double or guardian spirit who is both molded by and inspires the individual, leading them towards heaven. The magician (as Frabato knew) required symbols, images, ideas toward which to strive and focus. What better \emph{symbolon} than that which is beautiful, true, and good? Did not Saint Paul command meditation upon such? Where is this to be found, in the modern world, except in the past? The material may be around us, but the form lies in the past.

There is no necessary conflict between esoteric Christianity, alchemy, and the paganism which animated the Mediterranean basin and flowered in the path of imperium to Rome. It is the magician who will bring it to pass.


\flrightit{Posted on 2011-10-08 by Logres}

\begin{center}* * *\end{center}

\begin{footnotesize}
\begin{sffamily}

\texttt{logres on 2011-10-12 at 20:28 said:}

My posts are often of a rhetorical-poetic nature, although I struggle against this -- Cologero has more to offer for substance (I recommend the 400+ articles he has penned -- pick one at random). I have, however, sought to learn from his substance, to incorporate such as I can, and to strenuously avoid a purely ``persuasive'' tone (such as one encounters in most political writing, including the Federalist Papers)

The concept of self-persuasion over ``suasion'' (I think this originates with Weineger(?) and another German who influenced Evola) is critical to understanding what it means to separate from the modern world, in my humble opinion. This is partly ``additive'' to Guenon, etc. Sort of a ``negative'' morality lesson, if you will.

\hfill

\texttt{logres on 2011-10-12 at 20:30 said:}

It is on this basis of self-suasion that pagan and Christian can communicate.

\hfill

\texttt{logres on 2011-10-12 at 20:37 said:}

Here is a good one: \url{http://www.gornahoor.net/?p=3077\#more-3077}

I try to ask myself why I am asking other people questions, since every question to someone else is a question to the Self. To solve that puzzle, would be to answer your question immediately. You would see what is to be done. You already know the answer; if I give it to you, I have to wear a mask.

\hfill

\texttt{Charlotte Cowell on 2011-10-18 at 18:19 said:}

This is a wonderful post thank you -- my previous comment on initiation would have been better suited to this page. A thought with respect to `self initiation'. Perhaps there is (an understandable) blurring of distinction between being alone in one's work of the soul, one's own `Great Work' -- the pathworking of initiation that one walks alone for so much of the way -- and the initialising event, the catalyst (or catalysts, plural)? An initiated being will indeed be left very much `alone' to wrestle with his/her angel (and demons), yet will also (for great portions of time), be fully conscious of membership of a school -- even successive schools, as in primary, secondary, college, university -- and also of kindred spirits, fellow hermeticists shining their lamps if you will. With respect to rebirths, the soul/spirit in question must of course be `ripe' for the picking, so in that respect they do not so much self-initiate as they reach a point in their evolution where they're able to move beyond the self into contact with the divine hierarchy, to be conscious of the Holy Spirit, receptive to the supernatural. In this state of consciousness one cannot really be alone -- knowing from the onset that separation is a form of illusion -- although surely one might feel extremely lonely for long periods of time as the essential self-work is performed. Crossing the abyss is not likely to feel blissfully happy for the average being; even Christ descended to hell for a time... 

\hfill

\end{sffamily}
\end{footnotesize}

